\section{Common measurement noise and the fifth-power loss of information}
\label{sec:detectors}

This section derives consequences of the homogeneous-control theorem. These
are mathematical measurement models, not experimental validation. The phase
variance is $s>0$, the number of sites is $n=63$, and
\[
 X_s=I_A(s),\qquad Y_s=I_B(s),\qquad 0\le X_s,Y_s\le n^2.
\]
Their first three moments agree exactly. Write
\[
 D_j(s)=\mathbb E X_s^j-\mathbb E Y_s^j,\qquad
 C=\frac{455015761}{20000},\qquad c_5=16C=\frac{455015761}{1250}.
\]
The exact fourth-moment identity gives
\[
 D_4(s)=C e^{-4s}(1-e^{-s})^5
 (e^{-3s}+3e^{-2s}+6e^{-s}+6)>0,
 \qquad D_4(s)=c_5s^5+O(s^6).
\]
The separately computed fifth-moment polynomial satisfies
$D_5(s)=O(s^6)$. This also follows analytically: the scaled intensities
$I_G(s)/s$ tend in every fixed moment to the squared modulus of a complex
Gaussian with common covariance and common pseudocovariance magnitude;
therefore the $s^5$ coefficients of the fifth moments agree. Every moment
is a finite exponential polynomial, so its Taylor expansion has integer
powers. The saved rational calculation verifies all coefficients through
order five of $D_5$ are zero without using this limiting argument.

\subsection{A uniform moment bound}

Let $Z\sim N(0,I_n)$, $P=I-\mathbf1\mathbf1^T/n$, and $U=PZ$.
Centering changes the amplitude by a common phase and leaves intensity
unchanged. At the common clean null, $\sum_i z_i=0$, so
\[
 \left|\sum_i z_i e^{i\sqrt{s}U_i}\right|
 =\left|\sum_i z_i(e^{i\sqrt{s}U_i}-1)\right|
 \le\sqrt{ns}\,\|U\|.
\]
Since $\|U\|^2$ is chi-square with $n-1$ degrees of freedom,
\begin{equation}\label{eq:Bp}
 \mathbb E X_s^p+\mathbb E Y_s^p\le B_p s^p,
 \qquad B_p=2n^p\prod_{j=0}^{p-1}(n-1+2j).
\end{equation}
The identity for the chi-square moment follows by differentiating its
moment generating function $(1-2t)^{-(n-1)/2}$ at zero. This estimate is
deliberately conservative and holds for every $s>0$.

\subsection{Distinction survives common additive noise}

Let $E$ have any fixed distribution with finite fourth absolute moment,
independent of the intensity, and use the same noise law for A and B.
The binomial expansion proves that the first three moments of $X_s+E$
and $Y_s+E$ agree and
\[
 \mathbb E(X_s+E)^4-\mathbb E(Y_s+E)^4=D_4(s)>0.
\]
Thus the measured distributions remain distinct at every finite positive
$s$. This is an identifiability statement; it supplies no uniform practical
accuracy guarantee as the noise level grows or the signal shrinks.

\subsection{Gaussian readout: an exact leading total-variation asymptotic}

Take $E\sim N(0,\sigma^2)$ with fixed $\sigma>0$, and let $g_\sigma$ be its
density. Denote total variation by $d_{\rm TV}(P,Q)=\tfrac12\|p-q\|_1$
when densities exist. Taylor's theorem with an integral remainder gives,
in $L^1(\mathbb R)$,
\[
 p_A-p_B=\frac{D_4(s)}{24}g_\sigma^{(4)}
             -\frac{D_5(s)}{120}g_\sigma^{(5)}+R_s,
 \qquad
 \|R_s\|_1\le\frac{B_6s^6}{720}\|g_\sigma^{(6)}\|_1.
\]
All terms of degree zero through three cancel. Translation invariance of
the $L^1$ norm justifies the remainder bound for all nonnegative input
values, including unbounded Taylor segments. Consequently,
\begin{equation}\label{eq:gaussian-asymptotic}
 d_{\rm TV}(X_s+E,Y_s+E)
 =\frac{c_5}{48}\|g_\sigma^{(4)}\|_1s^5+O(s^6).
\end{equation}
The remainder may depend on the fixed detector parameter $\sigma$.
In particular the distance is of order $s^5$, with a strictly positive
coefficient; this statement concerns total variation, not merely one
selected statistic.

For explicit constants, let $H_k$ be the probabilists' Hermite polynomial.
Then
\[
 \|g_\sigma^{(k)}\|_1
 =\sigma^{-k}\mathbb E|H_k(Z)|\le\frac{\sqrt{k!}}{\sigma^k},
 \qquad Z\sim N(0,1).
\]
The inequality follows from Cauchy--Schwarz and
$\mathbb E H_k(Z)^2=k!$, which follows by repeated Gaussian integration
by parts. Put $a=\sqrt{3-\sqrt6}$, $b=\sqrt{3+\sqrt6}$, and
$\varphi(t)=e^{-t^2/2}/\sqrt{2\pi}$. Splitting at the four roots of
$H_4(t)=t^4-6t^2+3$ and integrating
$(H_3\varphi)'=-H_4\varphi$ yields
\[
 \|g_\sigma^{(4)}\|_1
 =\frac{4\sqrt6}{\sigma^4}\bigl[a\varphi(a)+b\varphi(b)\bigr]>0.
\]

The saved arithmetic supplies rational $A_4,A_5$ such that
\[
 |D_4(s)-c_5s^5|\le A_4s^6,
 \qquad |D_5(s)|\le A_5s^6.
\]
For a polynomial $D_j(s)=\sum_m a_{jm}e^{-ms}$ they are obtained from
$A_j=\sum_m|a_{jm}|m^6/6!$ after the exact lower-order cancellations.
Thus an explicit finite-$s$ error in \eqref{eq:gaussian-asymptotic} is
\begin{equation}\label{eq:gaussian-error}
 s^6\left[
 \frac{A_4\|g_\sigma^{(4)}\|_1}{48}
 +\frac{A_5\|g_\sigma^{(5)}\|_1}{240}
 +\frac{B_6\|g_\sigma^{(6)}\|_1}{1440}\right].
\end{equation}
A simpler all-$s$ upper bound, using a Taylor expansion only through
degree four, is
\[
 d_{\rm TV}(X_s+E,Y_s+E)
 \le\min\left\{1,
 \frac{D_4(s)}{2\sqrt{24}\,\sigma^4}
 +\frac{B_5s^5}{2\sqrt{120}\,\sigma^5}\right\}.
\]

\subsection{Photon counting: strict distinction and fifth-power scaling}

Consider a count $N$ conditional on intensity $I$ with
\[
 N\mid I\sim\operatorname{Poisson}(b+\eta I),
 \qquad b\ge0,\quad\eta>0,
\]
where the background $b$ and gain $\eta$ are fixed and common to the two
hypotheses. The falling factorial $(N)_j=N(N-1)\cdots(N-j+1)$ satisfies
$\mathbb E[(N)_j\mid I]=(b+\eta I)^j$. Hence the first three factorial
moments agree and the fourth differs by $\eta^4D_4(s)>0$. The two count
distributions are therefore distinct for every finite $s>0$.

Let $p_b$ be the Poisson mass function at mean $b$, and let the shift
operator $S$ act by $(Sp)(k)=p(k-1)$ with $p(-1)=0$. Derivatives with
respect to intensity satisfy
\[
 \frac{d^j}{dI^j}p_{b+\eta I}
 =\eta^j(S-\mathrm{Id})^j p_{b+\eta I},
 \qquad \left\|\frac{d^j}{dI^j}p_{b+\eta I}\right\|_1\le(2\eta)^j.
\]
Taylor expansion in the Banach space $\ell^1$ therefore proves
\begin{equation}\label{eq:poisson-asymptotic}
 d_{\rm TV}(N_A,N_B)
 =\frac{c_5\eta^4}{48}\|(S-\mathrm{Id})^4p_b\|_1s^5+O(s^6).
\end{equation}
The leading constant is positive: the generating function of
$(S-\mathrm{Id})^4p_b$ is $(z-1)^4e^{b(z-1)}$, which is not the zero
function. For zero background, $p_0$ is a point mass and the fourth
difference has coefficients $(1,-4,6,-4,1)$, so
\begin{equation}\label{eq:poisson-zero}
 d_{\rm TV}(N_A,N_B)=\frac{c_5\eta^4}{3}s^5+O(s^6).
\end{equation}
An explicit all-$s$ error for this last expression is
\[
 s^6\left[\frac{\eta^4A_4}{3}
 +\frac{2\eta^5A_5}{15}+\frac{2\eta^6B_6}{45}\right].
\]
An all-$s$ upper bound requiring no sixth-order quantities is
\[
 d_{\rm TV}(N_A,N_B)\le\min\left\{1,
 \frac{\eta^4D_4(s)}{3}+\frac{2\eta^5B_5s^5}{15}\right\}.
\]

\subsection{An asymptotically optimal count-parity rule}

In the zero-background model, the difference in even-count probabilities is
\[
 \Pr(N_A\text{ even})-\Pr(N_B\text{ even})
 =\frac12\left(\mathbb E e^{-2\eta X_s}
                   -\mathbb E e^{-2\eta Y_s}\right)
 =\frac{c_5\eta^4}{3}s^5+O(s^6).
\]
The same coefficient occurs in \eqref{eq:poisson-zero}. With equal prior
probabilities and one observation, choosing A on even counts and B on odd
counts therefore has advantage over random guessing asymptotic to the
optimal advantage: the ratio tends to one as $s\downarrow0$.
This is a one-observation asymptotic statement. Its absolute advantage
tends to zero, and the argument does not claim finite-$s$ optimality,
optimality for positive background, or optimality for multiple observations.

\subsection{Sampling and interpretation}

For $m$ independent observations from one unknown configuration with equal
priors, a one-observation TV bound $\delta(s)$ implies
\[
 \Pr(\text{correct})\le\frac{1+\min\{1,m\delta(s)\}}2.
\]
This elementary inequality gives a necessary $\Omega(s^{-5})$ sample bound.
Section~\ref{sec:sampling} improves it to sharp necessary and sufficient
orders using Hellinger distance. Unknown gain, background, noise calibration,
temporal dependence, and changes in the structural model are outside the
stated experiment.

